% =============================================================================
% Chapter 5  The Standard Form, Encoded Exactly  --  ~3.5 pp
% rubric: Project Body (Modelling) -- Model Formulation /30.
% Plan: ThesisPlanning.md §5.5 (supporting contribution S2).
%
% PURPOSE     Answer question (i) of sec:pihm-assessment: build the paper's own
%             encoding as published, then the best EXACT encoding of the paper's
%             regime, and show that even the best encoder leaves a Theta~(N^4)
%             degree -- the cost is in the HAMILTONIAN, not the encoder.
% CRITERION   Formulation band 4, "providing a significant advance in the state
%             of the art": S2 is a genuine, verified construction. Frame it as
%             serving the comparison (a fairness device), never as a rival solver.
% SOURCES     Planning.md §3.7 (the published build; the structured exact
%             encoding: factorisation, Lambda, U_odd, R_0, lifts, rank-one rows,
%             products with ancilla reuse), §6.3; App. A-4, A-6; claims A-06 to
%             A-08, A-12, C-01, C-02; pihm/src/pihm/circuits/standard.py,
%             arithmetic.py; pihm/docs/PILLAR1.md.
% WRITE WHEN  Now (wave W1; P3 done); §5.4's numbers at gate G3 (wave W2).
% =============================================================================

\subsection{The Published Encoding, Rebuilt}
\label{sec:standard-published}
% ~0.75 pp
% LAND, in order:
%   1. Figs 2, 8 and 9 of the paper built as drawn, from the legible figure and
%      the caption angles (sec:pihm-encoding describes them).
%   2. What is measured, claim by claim: qubits against the claimed 2n + 3;
%      gates after multi-controlled-X decomposition against the claimed O(n^3);
%      the entry error eps_G(n) against "exponentially decreasing" and against
%      the small-angle bound (prop:epsilon-G, Appendix D); the prefactor
%      (eq:pihm-published-prefactor).
%   3. The results in one sentence, pointing to sec:results-operators. THE key
%      number: alpha / ||G||_2 = 2.2e6 at n = 7, growing like N^3 (\prov).
% OFFLOAD: -> Appendix D (app:pihm-circuit): "every deviation from the printed
%   figure, and the measured qubits, gates and eps_G(n), are given in
%   Appendix D".

\subsection{An Exact Factorisation of the Differentiation Matrix}
\label{sec:standard-factorisation}
% ~1 pp
% LAND, in order:
%   1. The idea in words FIRST: every entry of G is (a row-0 weight) x (a
%      parity-selected shift) x (the column index).
%   2. prop:G-factorisation (claim C-01). The proposition is exact; quote the
%      measured 2.3e-13 (n <= 10) as VERIFICATION, not as evidence.
% OFFLOAD: proof -> Appendix E (app:proof-G-factorisation).

\begin{proposition}[Exact factorisation of the differentiation matrix]
\label{prop:G-factorisation}
% STATE: for every n >= 1 the normalised-basis differentiation matrix
%   factorises exactly as below, with R_0 a row-0 reweighting, U_odd the
%   average of the N/2 truncated shifts by odd offsets (times N/2) and Lambda
%   the column index.
\begin{equation}
    \begin{gathered}
        \G = R_0\,\bigl(2\Uodd\bigr)\,\Lambda, \\
        R_0 = I - \Bigl(1 - \tfrac{1}{\sqrt2}\Bigr)\ket{0}\bra{0},
        \qquad
        \bigl(\Uodd\bigr)_{kp} = \bigl[\,p > k,\ p - k\ \text{odd}\,\bigr],
        \qquad
        \Lambda = \operatorname{diag}(p) .
    \end{gathered}
    \label{eq:G-structured}
\end{equation}
\end{proposition}

\subsection{Circuits and Cost}
\label{sec:standard-cost}
% ~1 pp
% LAND, in order:
%   1. Lambda as a Z-string LCU (eq:lambda-lcu): alpha = N - 1,
%      ceil(log2(n+1)) selector qubits, O(n) gates.
%   2. U_odd: a uniform PREPARE over odd offsets (X on the offset register's
%      lowest bit, H on the others); SELECT as the in-place subtraction
%      p <- p - d with a post-selected borrow flag, so wrapped terms vanish
%      EXACTLY (alpha = N/2, O(n) Toffolis). This is the one representative
%      circuit (F3).
%   3. R_0: a two-term LCU of I and I - 2|0><0|, alpha = 1.
%   4. prop:standard-cost. The Omega(n) ancilla floor comes from the N/2-term
%      LCU -- a genuine regime difference from the descriptor's O(log n).
%   5. Lifts M_{x^p} (banded shift x diagonal LCUs, shared with the descriptor);
%      rank-one rows via the feature map; products with ancilla reuse.
% OFFLOAD: the arithmetic SELECT and the resource derivation -> Appendix E
%   (app:structured).
% FIG: fig:uodd-select (F3). TAB (optional): tab:alpha-G-short; the full A-4
%   table is tab:alpha-G-full in Appendix E.

\begin{equation}
    \Lambda = \frac{N-1}{2}\,I - \sum_{j=0}^{n-1} 2^{j-1} Z_j .
    \label{eq:lambda-lcu}
\end{equation}

\placeholderfigure{SELECT for the odd-offset shifts: in-place subtraction with a borrow flag.}{fig:uodd-select}{Quantikz circuit: offset register (PREPARE: $X$ on the lowest bit, $H$ on the rest), system register $p$, borrow-flag ancilla; a ripple-carry subtractor $p \leftarrow p - d$ \cite{cuccaro2004adder} whose borrow flag is post-selected on $\ket{0}$, so every wrapped term vanishes exactly.}

\begin{proposition}[Cost of the structured encoding]
\label{prop:standard-cost}
% STATE: the factorisation of prop:G-factorisation compiles to an exact block
%   encoding of G with the subnormalisation, gate count and ancilla count below.
\begin{equation}
    \alpha_{\G} = N(N-1),
    \qquad
    \frac{\alpha_{\G}}{\|\G\|_2} \le 2.12 \quad (n \le 10),
    \qquad
    O(n^2)\ \text{gates},
    \qquad
    n + O(\log n)\ \text{ancilla} .
    \label{eq:standard-cost}
\end{equation}
\end{proposition}
% REVISE when writing: "<= 2.12 (n <= 10)" is a measured bound (A-4, 1.62 at
%   n = 2 rising to 2.115 at n = 10), not part of what the proposition proves --
%   state it beside the proposition, quoted with
%   \res[4]{derivative/n10}{structured_ratio}.

\begin{table}[htbp]
    \centering
    \caption{Subnormalisation of the differentiation-matrix encodings, relative to $\|\G\|_2$.}
    \label{tab:alpha-G-short}
    \small
    \begin{tabular}{@{}rrrr@{}}
        \toprule
        $n$ & $\|\G\|_2$ & Structured (\Cref{prop:standard-cost}) & Published (\Cref{eq:pihm-published-prefactor}) \\
        \midrule
        2 & \res[3,fixed]{derivative/n2}{norm} & \res[4]{derivative/n2}{structured_ratio} & \res[3]{derivative/n2}{published_ratio} \\
        4 & \res[3,fixed]{derivative/n4}{norm} & \res[4]{derivative/n4}{structured_ratio} & \res[3,sci]{derivative/n4}{published_ratio} \\
        7 & \res[3,fixed]{derivative/n7}{norm} & \res[4]{derivative/n7}{structured_ratio} & \res[3,sci]{derivative/n7}{published_ratio} \\
        10 & \res[3,fixed]{derivative/n10}{norm} & \res[4]{derivative/n10}{structured_ratio} & \res[3,sci]{derivative/n10}{published_ratio} \\
        \bottomrule
    \end{tabular}
\end{table}

\subsection{What Remains: The Standard-Form Hamiltonian}
\label{sec:standard-hamiltonian}
% ~0.75 pp
% LAND, in order:
%   1. alpha_H / ||H|| stays bounded (22 -> 60 over n = 2-8, Fig. 3b; \prov):
%      the encoder has done its job (claim C-02).
%   2. But ||H_std|| = Theta(N^8) against a flat gap (7.04), so by eq:degree the
%      degree is Theta~(N^4): 1.2e10 at n = 8 (\prov) -- eq:standard-wall.
%   3. THE HAMILTONIAN WALL, interpreted: ||G|| = Theta(N^2), squared by the Gram
%      product and again by second derivatives; no encoder can undo a norm.
%   4. Close by naming Chapter 6 (sec:descriptor) and what it changes: the
%      Hamiltonian itself.

\begin{equation}
    \|\Hstd\| = \Theta(N^8),
    \qquad
    \Delta = \Theta(1)
    \quad\Longrightarrow\quad
    d = \tilde\Theta\bigl(N^4\bigr) .
    \label{eq:standard-wall}
\end{equation}

% CHECKLIST (marker)
%   1. Is the published build reported as a TEST of the paper's claims, with a
%      verdict for each?
%   2. Is the factorisation a proposition, with its proof pointed to by content?
%   3. Does §5.4 make the encoder-versus-Hamiltonian distinction unmissable?
%   4. Is S2 framed as serving the comparison?
% TRAPS  Calling S2 "our standard-form solver"; quoting eps_G bounds that belong
%        to the legacy reconstruction rather than to the Fig. 9 build.
